\documentclass[12pt,a4paper]{article}

\usepackage{fontawesome5}
\providecommand{\faPencilAlt}{\faIcon{pencil-alt}}

\usepackage{styles/couleurs}
\usepackage{styles/mise_en_page}
\usepackage{styles/boites}
\usepackage{styles/maths}
\usepackage{styles/figures}

\usepackage{enumitem}
\usepackage{array}
\usepackage{tabularx}
\usepackage{amsmath}
\usepackage{amssymb}
\usepackage{graphicx}
\usepackage{tikz}
\usepackage{xparse}

\usetikzlibrary{calc,angles,quotes,arrows.meta}

\setlist[itemize]{leftmargin=1.4em,itemsep=0.15em,topsep=0.2em}
\setlist[enumerate]{leftmargin=*,itemsep=0.22em,topsep=0.25em}

\newcommand{\lignepointillee}{\noindent\dotfill\par\vspace{0.45em}}

\newcommand{\lignesreponse}[1]{%
\begin{tcolorbox}[
enhanced,colback=white,colframe=NeutreC!40,boxrule=0.5pt,
arc=1.5mm,left=3mm,right=3mm,top=2mm,bottom=1mm,
before skip=2pt,after skip=2pt]
\foreach \i in {1,...,#1}{\lignepointillee}
\end{tcolorbox}}

\newcommand{\titreseq}[1]{%
\begin{center}
{\sffamily\large\bfseries\color{SavoirC}#1}
\end{center}}

\begin{document}

\titreseq{Automatismes -- Séries courtes}

\section*{Série 1 -- Triangle rectangle}

Dans le triangle \(ABC\) rectangle en \(B\) :

\begin{enumerate}
  \item nomme l'hypoténuse ;
  \item nomme les deux angles aigus ;
  \item nomme le côté adjacent à \(\widehat A\) ;
  \item nomme le côté adjacent à \(\widehat C\).
\end{enumerate}

\section*{Série 2 -- Écrire un cosinus}

Dans le triangle \(DEF\) rectangle en \(E\), complète :

\[
\cos(\widehat D)=\ldots\div\ldots
\]

\[
\cos(\widehat F)=\ldots\div\ldots
\]

\section*{Série 3 -- Calculatrice}

Donne une valeur approchée au centième.

\[
\cos 20^\circ
\qquad
\cos 35^\circ
\qquad
\cos 60^\circ
\qquad
\cos 75^\circ
\]

Quel calcul permet de contrôler rapidement que la calculatrice est en mode degré ?

\section*{Série 4 -- Multiplier ou diviser ?}

Pour chaque cas, indique seulement l'opération.

\begin{enumerate}
  \item Hypoténuse \(=12\) cm, angle \(=40^\circ\), côté adjacent cherché.
  \item Côté adjacent \(=7\) cm, angle \(=52^\circ\), hypoténuse cherchée.
  \item Hypoténuse \(=5{,}4\) m, angle \(=18^\circ\), côté adjacent cherché.
\end{enumerate}

\section*{Série 5 -- Calcul direct}

Calcule au dixième.

\[
8\times\cos 37^\circ
\qquad
15\times\cos 64^\circ
\]

\[
6\div\cos 42^\circ
\qquad
3{,}5\div\cos 70^\circ
\]

\section*{Série 6 -- Compléter une rédaction}

Dans le triangle \(GHI\) rectangle en \(H\), on connaît
\(GI=9\) cm et \(\widehat G=31^\circ\).

\[
\cos(\widehat G)=\ldots\div\ldots
\]

\[
GH=\ldots\times\cos\ldots
\]

\[
GH\approx\ldots\text{ cm}.
\]

\section*{Série 7 -- Vrai ou faux ?}

\begin{enumerate}
  \item Le cosinus d'un angle aigu est une longueur.
  \item Le côté adjacent est toujours plus petit que l'hypoténuse.
  \item Pour calculer une hypoténuse, on multiplie le côté adjacent par le cosinus.
  \item Si l'angle augmente, son cosinus diminue.
\end{enumerate}

\section*{Série 8 -- Repérer une erreur}

Un élève écrit :

\[
\text{Dans }ABC\text{ rectangle en }B,\quad
AC=6\times\cos 40^\circ=4{,}6\text{ cm}.
\]

On sait que \(AB=6\) cm et que \(AC\) est l'hypoténuse.

Explique l'erreur et corrige le calcul.

\end{document}
